% GENERATED FILE. Edit canonical lessons, then run npm run generate:books.
% canonical-source-hash: fnv1a64:fcd173f64dcef045
% canonical-lessons: ML-C342-pokanam, ML-C342-audyogikamaya, ML-W342-au, ML-C342-vannatukondu, ML-R342-advice-and-reasons

\chapter{Must, Ought To, and Because}
\label{ch:ml-must-ought-to-and-because}
\hlchaptermodality{342}

\begin{chapteropening}
\textbf{By the end of this chapter:} \emph{I can say what I must, ought to and need not do, and give a reason for an action or a plan with -atukoṇṭŭ or kāraṇaṁ.}

Complete ML-R342-advice-and-reasons: must, ought to and need not, reasons with -atukoṇṭŭ and kāraṇaṁ, a plan with its reason, and the new letter \ml{ഔ}.
\end{chapteropening}

\section[\texorpdfstring{\ml{ഞാൻ} \ml{പോകണം} (ñān pōkaṇaṁ)}{ഞാൻ പോകണം (ñān pōkaṇaṁ)}]{\ml{ഞാൻ} \ml{പോകണം} (ñān pōkaṇaṁ) — I must go}

\label{lesson:ML-C342-pokanam}



\begin{quote}
Say \textquotedblleft{}I want tea\textquotedblright{}. (\emph{enikkŭ cāya vēṇaṁ}.) Then \textquotedblleft{}no need, I don't want it\textquotedblright{}. (\emph{vēṇṭa}.)
\end{quote}

\begin{grammarlens}[title={Grammar Lens: -\ml{അണം} on the verb}]
Take a verb's dictionary form, drop \textbf{-\ml{ഉക}}, and add \textbf{-\ml{അണം}}:

\begin{quote}
\ml{പോകുക} $\to$ \textbf{\ml{പോകണം}}, \emph{pōkaṇaṁ} · \ml{വരുക} $\to$ \textbf{\ml{വരണം}}, \emph{varaṇaṁ} · \ml{പഠിക്കുക} $\to$ \textbf{\ml{പഠിക്കണം}}, \emph{paṭhikkaṇaṁ}
\end{quote}

Grammars explain \textbf{-\ml{അണം}} as \textbf{\ml{വേണം}} itself, worn down and joined to the verb. Who stands in front decides what it means:

\begin{quote}
\textbf{\ml{ഞാൻ} \ml{നാളെ} \ml{പോകണം}.}
\end{quote}

\begin{quote}
\emph{ñān nāḷe pōkaṇaṁ.} — \textquotedblleft{}I must go tomorrow.\textquotedblright{}
\end{quote}

\begin{quote}
\textbf{\ml{എനിക്ക്} \ml{നാളെ} \ml{പോകണം}.}
\end{quote}

\begin{quote}
\emph{enikkŭ nāḷe pōkaṇaṁ.} — \textquotedblleft{}I want to go tomorrow.\textquotedblright{}
\end{quote}

With \textbf{\ml{ഞാൻ}} it is a duty; with the dative \textbf{\ml{എനിക്ക്}}, as with \textbf{\ml{വേണം}} alone, it is a wish or a need.
\end{grammarlens}

\begin{grammarlens}[title={Grammar Lens: ought to, need not, must not}]
\textbf{-\ml{ഏണ്ടതാണ്}} says what is right to do: \textquotedblleft{}ought to\textquotedblright{}.

\begin{quote}
\textbf{\ml{നീ} \ml{നേരത്തെ} \ml{ഉറങ്ങേണ്ടതാണ്}.}
\end{quote}

\begin{quote}
\emph{nī nēratte uṟaṅṅēṇṭatāṇŭ.} — \textquotedblleft{}You ought to sleep early.\textquotedblright{}
\end{quote}

\textbf{-\ml{ഏണ്ട}}, the \textbf{\ml{വേണ്ട}} you know, says \textquotedblleft{}need not\textquotedblright{}: \textbf{\ml{നീ} \ml{വരേണ്ട}} (\emph{nī varēṇṭa}), \textquotedblleft{}you need not come\textquotedblright{}. \textbf{-\ml{അരുത്}} forbids: \textbf{\ml{നീ} \ml{വൈകരുത്}} (\emph{nī vaikarutŭ}), \textquotedblleft{}you must not be late\textquotedblright{}.
\end{grammarlens}

\subsection*{Your turn}
Say these aloud:

\begin{itemize}
\raggedright
  \item \textquotedblleft{}I must go tomorrow\textquotedblright{}, then \textquotedblleft{}I want to go tomorrow\textquotedblright{}
  \item \textquotedblleft{}you must come\textquotedblright{}, \textquotedblleft{}I must study\textquotedblright{} — \emph{nī varaṇaṁ}, \emph{ñān paṭhikkaṇaṁ}
  \item \textquotedblleft{}You ought to sleep early.\textquotedblright{}
  \item \textquotedblleft{}You need not come\textquotedblright{}, then \textquotedblleft{}you must not be late\textquotedblright{}
\end{itemize}

\begin{tcolorbox}[breakable,colback=teal!4,colframe=teal!35!black,title={Before you move on}]
\textquotedblleft{}I must go\textquotedblright{}? (\emph{\textbf{ñān pōkaṇaṁ.}}) What does \emph{enikkŭ pōkaṇaṁ} say instead? (\textbf{I want to go.}) \textquotedblleft{}Ought to sleep\textquotedblright{}? (\emph{\textbf{uṟaṅṅēṇṭatāṇŭ.}}) \textquotedblleft{}Need not come\textquotedblright{}? (\emph{\textbf{varēṇṭa.}})
\end{tcolorbox}
\section[\texorpdfstring{\ml{ഔദ്യോഗികമായ} (audyōgikamāya)}{ഔദ്യോഗികമായ (audyōgikamāya)}]{\ml{ഔദ്യോഗികമായ} (audyōgikamāya) — official}

\label{lesson:ML-C342-audyogikamaya}



\begin{quote}
Before the new one: say \textquotedblleft{}an official, an officer\textquotedblright{}, then \textquotedblleft{}a document\textquotedblright{}. (\emph{udyōgasthan}, \emph{rēkha}.) Then \textquotedblleft{}I must go\textquotedblright{}. (\emph{ñān pōkaṇaṁ}.)
\end{quote}

\subsection*{You'll want to know: \ml{ഔദ്യോഗികമായ}}
\textbf{\ml{ഔദ്യോഗികമായ}} — \emph{audyōgikamāya} — \textquotedblleft{}official\textquotedblright{}: of an office, done formally.

\begin{quote}
\textbf{\ml{ഇത്} \ml{ഒരു} \ml{ഔദ്യോഗികമായ} \ml{രേഖയാണ്}.}
\end{quote}

\begin{quote}
\emph{itŭ oru audyōgikamāya rēkhayāṇŭ.} — \textquotedblleft{}This is an official document.\textquotedblright{}
\end{quote}

In set phrases it stands before its noun without \textbf{-\ml{മായ}}:

\begin{quote}
\textbf{\ml{മലയാളം} \ml{കേരളത്തിന്റെ} \ml{ഔദ്യോഗിക} \ml{ഭാഷയാണ്}.}
\end{quote}

\begin{quote}
\emph{malayāḷaṁ kēraḷattinṟe audyōgika bhāṣayāṇŭ.} — \textquotedblleft{}Malayalam is the official language of Kerala.\textquotedblright{}
\end{quote}

With \textbf{-\ml{ആയി}} it becomes \textquotedblleft{}officially\textquotedblright{}: \textbf{\ml{ഔദ്യോഗികമായി}} (\emph{audyōgikamāyi}). Its first letter is new; the next lesson takes it out and puts it in your hand.

\begin{cousinweb}
\textbf{\ml{ഔദ്യോഗികമായ}} (\emph{audyōgikamāya}) is a learned word built the Sanskrit way on \emph{udyoga}, \textquotedblleft{}undertaking, exertion\textquotedblright{} (Monier-Williams), which Malayalam took up for a post in an office: \textbf{\ml{ഉദ്യോഗസ്ഥൻ}}, \textquotedblleft{}an officer\textquotedblright{}, is the same root. Strengthening the first vowel, \emph{u} to \emph{au}, and adding \emph{-ika} makes \textquotedblleft{}having to do with\textquotedblright{} it, so \textquotedblleft{}of the office\textquotedblright{}. The Malayalam \textbf{-\ml{മായ}} then makes it an adjective.
\end{cousinweb}

\subsection*{Your turn}
Say these aloud:

\begin{itemize}
\raggedright
  \item \emph{audyōgikamāya}, in five beats: \emph{au-dyō-gi-ka-māya}
  \item \textquotedblleft{}This is an official document.\textquotedblright{}
  \item \textquotedblleft{}Malayalam is the official language of Kerala.\textquotedblright{}
  \item \textquotedblleft{}officially\textquotedblright{} — \emph{audyōgikamāyi}
\end{itemize}

\begin{tcolorbox}[breakable,colback=teal!4,colframe=teal!35!black,title={Before you move on}]
\textquotedblleft{}Official\textquotedblright{}? (\textbf{\ml{ഔദ്യോഗികമായ}}.) \textquotedblleft{}The official language\textquotedblright{}? (\emph{\textbf{audyōgika bhāṣa.}}) Which word for a person shares its root? (\textbf{\ml{ഉദ്യോഗസ്ഥൻ}}, \textquotedblleft{}an officer\textquotedblright{}.)
\end{tcolorbox}
\section[\texorpdfstring{\ml{ഔ} (au)}{ഔ (au)}]{\ml{ഔ} (au) — the letter that opens \ml{ഔദ്യോഗികമായ}}

\label{lesson:ML-W342-au}



\begin{quote}
Before the new one: write \textbf{\ml{ഓ}} and \textbf{\ml{ഊ}} once each, and name them.

Now say \textbf{\ml{ഔദ്യോഗികമായ}} (\emph{audyōgikamāya}), \textquotedblleft{}official\textquotedblright{}. This lesson takes its first letter out and puts it in your hand.
\end{quote}

\subsection*{Script you'll notice: \ml{ഔ}}
\textbf{\ml{ഔ}} — \emph{au}, at the start of a word. It is rare: a handful of learned words begin with it.

Its body is the shape of \textbf{\ml{ഒ}}, which you write inside \textbf{\ml{ഓ}}: an inner loop, a broad arch, a rounded lobe. On its right it carries the tail you have written on \textbf{\ml{ഊ}}: a small lobe, a notch, and a broader lobe that sweeps back along its foot. A later letter lesson is about that tail alone.

\subsection*{Writing: \ml{ഔ}}
\textbf{\ml{ഔ}} is two parts side by side: the body of \textbf{\ml{ഒ}} on the left, and on the right the tail you know from \textbf{\ml{ഊ}}. Write the body first, then add the tail beside it: a small lobe, a notch, and a broader lobe that sweeps back along its foot.

Write \textbf{\ml{ഔ}} yourself, slowly, and larger than it is printed. Copy the printed shape.

\subsection*{Your turn}
\begin{itemize}
\raggedright
  \item \emph{Find:} \ml{ഔ} at the start of \textbf{\ml{ഔദ്യോഗികമായ}}
  \item \emph{Trace it:} \ml{ഔ} three times, saying \emph{au} as you finish each one
  \item \emph{Write it:} \ml{ഊ}, then \ml{ഔ}, from memory, and compare their tails
  \item \emph{Write it:} \textbf{\ml{ഔദ്യോഗികമായ}}
\end{itemize}

\begin{tcolorbox}[breakable,colback=teal!4,colframe=teal!35!black,title={Before you move on}]
Which letter is this — \ml{ഔ}? (\textbf{au}, at the start of a word.) Which word did it come from? (\textbf{\ml{ഔദ്യോഗികമായ}}, \emph{audyōgikamāya}, \textquotedblleft{}official\textquotedblright{}.) Which other letter carries the same tail? (\textbf{\ml{ഊ}}.)
\end{tcolorbox}
\section[\texorpdfstring{\ml{വന്നതുകൊണ്ട്} (vannatukoṇṭŭ)}{വന്നതുകൊണ്ട് (vannatukoṇṭŭ)}]{\ml{വന്നതുകൊണ്ട്} (vannatukoṇṭŭ) — because he came}

\label{lesson:ML-C342-vannatukondu}



\begin{quote}
Say \textquotedblleft{}I had a fever. So I did not come.\textquotedblright{} Then the same with \emph{kāraṇaṁ}.
\end{quote}

\begin{grammarlens}[title={Grammar Lens: -\ml{അതുകൊണ്ട്} on a verb}]
To put the reason inside the sentence, turn its verb into \textquotedblleft{}the …-ing\textquotedblright{} with \textbf{-\ml{അത്}}, as in \emph{vaikiyatŭ}, and add \textbf{\ml{കൊണ്ട്}}, \textquotedblleft{}with\textquotedblright{}:

\begin{quote}
\ml{വന്നു} $\to$ \textbf{\ml{വന്നതുകൊണ്ട്}} — \emph{vannatukoṇṭŭ} — because (he) came
\end{quote}

\begin{quote}
\textbf{\ml{മഴ} \ml{പെയ്തതുകൊണ്ട്} \ml{ഞാൻ} \ml{വന്നില്ല}.}
\end{quote}

\begin{quote}
\emph{maḻa peytatukoṇṭŭ ñān vannilla.} — \textquotedblleft{}Because it rained, I did not come.\textquotedblright{}
\end{quote}

The present works the same way: \textbf{\ml{വരുന്നു}} gives \textbf{\ml{വരുന്നതുകൊണ്ട്}}. \textbf{\ml{ഇല്ല}} gives \textbf{\ml{ഇല്ലാത്തതുകൊണ്ട്}}:

\begin{quote}
\textbf{\ml{സമയം} \ml{ഇല്ലാത്തതുകൊണ്ട്} \ml{ഞാൻ} \ml{ബസിൽ} \ml{പോയി}.}
\end{quote}

\begin{quote}
\emph{samayaṁ illāttatukoṇṭŭ ñān basil pōyi.} — \textquotedblleft{}Because there was no time, I went by bus.\textquotedblright{}
\end{quote}
\end{grammarlens}

\begin{grammarlens}[title={Grammar Lens: a plan and its reason}]
A plan is in the future in \textbf{-\ml{ഉം}}. Give its reason after \textbf{\ml{കാരണം}}:

\begin{quote}
\textbf{\ml{അടുത്ത} \ml{വർഷം} \ml{ഞാൻ} \ml{കേരളത്തിൽ} \ml{പോകും}, \ml{കാരണം} \ml{എനിക്ക്} \ml{കൂടുതൽ} \ml{മലയാളം} \ml{പഠിക്കണം}.}
\end{quote}

\begin{quote}
\emph{aṭutta varṣaṁ ñān kēraḷattil pōkuṁ, kāraṇaṁ enikkŭ kūṭutal malayāḷaṁ paṭhikkaṇaṁ.} — \textquotedblleft{}Next year I will go to Kerala, because I want to learn more Malayalam.\textquotedblright{}
\end{quote}
\end{grammarlens}

\subsection*{Your turn}
Say these aloud:

\begin{itemize}
\raggedright
  \item \textquotedblleft{}Because it rained, I did not come.\textquotedblright{}
  \item \textquotedblleft{}Because there was no time, I went by bus.\textquotedblright{}
  \item next year's plan, and its reason after \emph{kāraṇaṁ}
  \item one thing you must do this week, and why
\end{itemize}

\begin{tcolorbox}[breakable,colback=teal!4,colframe=teal!35!black,title={Before you move on}]
\textquotedblleft{}Because he came\textquotedblright{}? (\emph{\textbf{vannatukoṇṭŭ.}}) \textquotedblleft{}Because there was no …\textquotedblright{}? (\emph{\textbf{illāttatukoṇṭŭ.}}) Which word gives the reason for a plan? (\textbf{\ml{കാരണം}}, after the plan.)
\end{tcolorbox}
\section[(review)]{Review — advice and reasons}

\label{lesson:ML-R342-advice-and-reasons}



\begin{quote}
Say \textquotedblleft{}official\textquotedblright{}. (\emph{audyōgikamāya}.) Write \textbf{\ml{ഔ}} once from memory.
\end{quote}

\subsection*{Your turn}
\begin{itemize}
\raggedright
  \item \emph{Say it:} \textquotedblleft{}I must go tomorrow\textquotedblright{}, then \textquotedblleft{}I want to go tomorrow\textquotedblright{}
  \item \emph{Say it:} \textquotedblleft{}You ought to sleep early. You must not be late.\textquotedblright{}
  \item \emph{Say it:} \textquotedblleft{}This is an official document.\textquotedblright{}
  \item \emph{Say it:} \textquotedblleft{}Because it rained, I did not come.\textquotedblright{}
  \item \emph{Say it:} next year's plan, with its reason after \emph{kāraṇaṁ}
  \item \emph{Say it:} \textquotedblleft{}I think you must come. So I am calling.\textquotedblright{}
  \item \emph{Write it:} \textbf{\ml{ഔദ്യോഗിക} \ml{ഭാഷ}}, then \textbf{\ml{ഐക്യം}}
\end{itemize}

\begin{tcolorbox}[breakable,colback=teal!4,colframe=teal!35!black,title={Before you move on}]
\textquotedblleft{}I must\textquotedblright{}? (\emph{\textbf{ñān pōkaṇaṁ}}: \textbf{-\ml{അണം}} on the verb.) \textquotedblleft{}Ought to\textquotedblright{}? (\textbf{-\ml{ഏണ്ടതാണ്}}.) \textquotedblleft{}Because it rained\textquotedblright{}? (\emph{\textbf{peytatukoṇṭŭ.}}) \textquotedblleft{}Official\textquotedblright{}? (\textbf{\ml{ഔദ്യോഗികമായ}}, which opens with \textbf{\ml{ഔ}}.)
\end{tcolorbox}
